\documentclass[10pt,a4paper]{amsart}
\usepackage{amsmath,amssymb,amsthm,graphicx}
% For commutative diagrams we use Xy-pic
\usepackage[all]{xy}

% We use 2cell for 2-commutative diagrams.
\xyoption{2cell}
% Theorem environments.
%
\theoremstyle{plain}
\newtheorem{theorem}[subsection]{Theorem}
\newtheorem{proposition}[subsection]{Proposition}
\newtheorem{lemma}[subsection]{Lemma}

\theoremstyle{definition}
\newtheorem{definition}[subsection]{Definition}
\newtheorem{example}[subsection]{Example}
\newtheorem{exercise}[subsection]{Exercise}
\newtheorem{situation}[subsection]{Situation}

\theoremstyle{remark}
\newtheorem{remark}[subsection]{Remark}
\newtheorem{remarks}[subsection]{Remarks}

\numberwithin{equation}{subsection}

% Macros
%
\def\lim{\mathop{\mathrm{lim}}\nolimits}
\def\colim{\mathop{\mathrm{colim}}\nolimits}
\def\Spec{\mathop{\mathrm{Spec}}}
\def\Hom{\mathop{\mathrm{Hom}}\nolimits}
\def\Ext{\mathop{\mathrm{Ext}}\nolimits}
\def\SheafHom{\mathop{\mathcal{H}\!\mathit{om}}\nolimits}
\def\SheafExt{\mathop{\mathcal{E}\!\mathit{xt}}\nolimits}
\def\Sch{\mathit{Sch}}
\def\Mor{\mathop{\mathrm{Mor}}\nolimits}
\def\Ob{\mathop{\mathrm{Ob}}\nolimits}
\def\Sh{\mathop{\mathit{Sh}}\nolimits}
\def\NL{\mathop{N\!L}\nolimits}
\def\CH{\mathop{\mathrm{CH}}\nolimits}
\def\proetale{{pro\text{-}\acute{e}tale}}
\def\etale{{\acute{e}tale}}
\def\QCoh{\mathit{QCoh}}
\def\Ker{\mathop{\mathrm{Ker}}}
\def\Im{\mathop{\mathrm{Im}}}
\def\Coker{\mathop{\mathrm{Coker}}}
\def\Coim{\mathop{\mathrm{Coim}}}

% Boxtimes
%
\DeclareMathSymbol{\boxtimes}{\mathbin}{AMSa}{"02}

%
% Macros for moduli stacks/spaces
%
\def\QCohstack{\mathcal{QC}\!\mathit{oh}}
\def\Cohstack{\mathcal{C}\!\mathit{oh}}
\def\Spacesstack{\mathcal{S}\!\mathit{paces}}
\def\Quotfunctor{\mathrm{Quot}}
\def\Hilbfunctor{\mathrm{Hilb}}
\def\Curvesstack{\mathcal{C}\!\mathit{urves}}
\def\Polarizedstack{\mathcal{P}\!\mathit{olarized}}
\def\Complexesstack{\mathcal{C}\!\mathit{omplexes}}
% \Pic is the operator that assigns to X its picard group, usage \Pic(X)
% \Picardstack_{X/B} denotes the Picard stack of X over B
% \Picardfunctor_{X/B} denotes the Picard functor of X over B
\def\Pic{\mathop{\mathrm{Pic}}\nolimits}
\def\Picardstack{\mathcal{P}\!\mathit{ic}}
\def\Picardfunctor{\mathrm{Pic}}
\def\Deformationcategory{\mathcal{D}\!\mathit{ef}}
\newenvironment{reference}{}{}
\newenvironment{slogan}{}{}
\newenvironment{history}{}{}
\usepackage{cleveref}
\crefname{originallemma}{original lemma}{original lemmas}
\Crefname{originallemma}{Original lemma}{Original lemmas}
\makeatletter
\newlabel{local-cohomology-lemma-finiteness-pushforwards-and-H1-local}{{0BK0}{source}{}{sourcepointer-1}{}}
\newlabel{local-cohomology-lemma-finiteness-pushforwards-and-H1-local@cref}{{[originallemma][0][]0BK0}{[1][0][]source}{}{}{}}
\makeatother
\crefname{originallemma}{original lemma}{original lemmas}
\Crefname{originallemma}{Original lemma}{Original lemmas}
\makeatletter
\newlabel{coherent-lemma-homs-over-open}{{01YB}{source}{}{sourcepointer-2}{}}
\newlabel{coherent-lemma-homs-over-open@cref}{{[originallemma][0][]01YB}{[1][0][]source}{}{}{}}
\makeatother
\crefname{originallemma}{original lemma}{original lemmas}
\Crefname{originallemma}{Original lemma}{Original lemmas}
\makeatletter
\newlabel{algebra-lemma-support-point}{{00L5}{source}{}{sourcepointer-3}{}}
\newlabel{algebra-lemma-support-point@cref}{{[originallemma][0][]00L5}{[1][0][]source}{}{}{}}
\makeatother
\crefname{originallemma}{original lemma}{original lemmas}
\Crefname{originallemma}{Original lemma}{Original lemmas}
\makeatletter
\newlabel{algebra-lemma-length-additive}{{00IV}{source}{}{sourcepointer-4}{}}
\newlabel{algebra-lemma-length-additive@cref}{{[originallemma][0][]00IV}{[1][0][]source}{}{}{}}
\makeatother
\crefname{originallemma}{original lemma}{original lemmas}
\Crefname{originallemma}{Original lemma}{Original lemmas}
\makeatletter
\newlabel{algebra-lemma-Artin-Rees}{{00IN}{source}{}{sourcepointer-5}{}}
\newlabel{algebra-lemma-Artin-Rees@cref}{{[originallemma][0][]00IN}{[1][0][]source}{}{}{}}
\makeatother
\crefname{originallemma}{original lemma}{original lemmas}
\Crefname{originallemma}{Original lemma}{Original lemmas}
\makeatletter
\newlabel{properties-lemma-quasi-affine-invertible-nonvanishing-section}{{0F20}{source}{}{sourcepointer-6}{}}
\newlabel{properties-lemma-quasi-affine-invertible-nonvanishing-section@cref}{{[originallemma][0][]0F20}{[1][0][]source}{}{}{}}
\makeatother
\crefname{originallemma}{original lemma}{original lemmas}
\Crefname{originallemma}{Original lemma}{Original lemmas}
\makeatletter
\newlabel{algebra-lemma-numerical-polynomial}{{00JZ}{source}{}{sourcepointer-7}{}}
\newlabel{algebra-lemma-numerical-polynomial@cref}{{[originallemma][0][]00JZ}{[1][0][]source}{}{}{}}
\makeatother
\crefname{originallemma}{original lemma}{original lemmas}
\Crefname{originallemma}{Original lemma}{Original lemmas}
\makeatletter
\newlabel{algebra-lemma-depth-drops-by-one}{{090R}{source}{}{sourcepointer-8}{}}
\newlabel{algebra-lemma-depth-drops-by-one@cref}{{[originallemma][0][]090R}{[1][0][]source}{}{}{}}
\makeatother
\crefname{originallemma}{original lemma}{original lemmas}
\Crefname{originallemma}{Original lemma}{Original lemmas}
\makeatletter
\newlabel{dualizing-lemma-depth}{{0AVZ}{source}{}{sourcepointer-9}{}}
\newlabel{dualizing-lemma-depth@cref}{{[originallemma][0][]0AVZ}{[1][0][]source}{}{}{}}
\makeatother
\crefname{originallemma}{original lemma}{original lemmas}
\Crefname{originallemma}{Original lemma}{Original lemmas}
\makeatletter
\newlabel{local-cohomology-lemma-finiteness-for-finite-locally-free}{{0BLT}{source}{}{sourcepointer-10}{}}
\newlabel{local-cohomology-lemma-finiteness-for-finite-locally-free@cref}{{[originallemma][0][]0BLT}{[1][0][]source}{}{}{}}
\makeatother
\crefname{originallemma}{original lemma}{original lemmas}
\Crefname{originallemma}{Original lemma}{Original lemmas}
\makeatletter
\newlabel{divisors-lemma-depth-pushforward}{{0EY0}{source}{}{sourcepointer-11}{}}
\newlabel{divisors-lemma-depth-pushforward@cref}{{[originallemma][0][]0EY0}{[1][0][]source}{}{}{}}
\makeatother
\crefname{originallemma}{original lemma}{original lemmas}
\Crefname{originallemma}{Original lemma}{Original lemmas}
\makeatletter
\newlabel{algebra-lemma-depth-in-ses}{{00LX}{source}{}{sourcepointer-12}{}}
\newlabel{algebra-lemma-depth-in-ses@cref}{{[originallemma][0][]00LX}{[1][0][]source}{}{}{}}
\makeatother
\crefname{originallemma}{original lemma}{original lemmas}
\Crefname{originallemma}{Original lemma}{Original lemmas}
\makeatletter
\newlabel{algebra-lemma-bound-depth}{{00LK}{source}{}{sourcepointer-13}{}}
\newlabel{algebra-lemma-bound-depth@cref}{{[originallemma][0][]00LK}{[1][0][]source}{}{}{}}
\makeatother
\crefname{originallemma}{original lemma}{original lemmas}
\Crefname{originallemma}{Original lemma}{Original lemmas}
\makeatletter
\newlabel{local-cohomology-lemma-local-finiteness-for-finite-locally-free}{{0BPY}{source}{}{sourcepointer-14}{}}
\newlabel{local-cohomology-lemma-local-finiteness-for-finite-locally-free@cref}{{[originallemma][0][]0BPY}{[1][0][]source}{}{}{}}
\makeatother
\crefname{originallemma}{original lemma}{original lemmas}
\Crefname{originallemma}{Original lemma}{Original lemmas}
\makeatletter
\newlabel{dualizing-lemma-special-case-local-duality}{{0AAK}{source}{}{sourcepointer-15}{}}
\newlabel{dualizing-lemma-special-case-local-duality@cref}{{[originallemma][0][]0AAK}{[1][0][]source}{}{}{}}
\makeatother
\crefname{originallemma}{original lemma}{original lemmas}
\Crefname{originallemma}{Original lemma}{Original lemmas}
\makeatletter
\newlabel{local-cohomology-lemma-sitting-in-degrees}{{0DWZ}{source}{}{sourcepointer-16}{}}
\newlabel{local-cohomology-lemma-sitting-in-degrees@cref}{{[originallemma][0][]0DWZ}{[1][0][]source}{}{}{}}
\makeatother
\crefname{originallemma}{original lemma}{original lemmas}
\Crefname{originallemma}{Original lemma}{Original lemmas}
\makeatletter
\newlabel{algebra-lemma-support-dimension-d}{{00L8}{source}{}{sourcepointer-17}{}}
\newlabel{algebra-lemma-support-dimension-d@cref}{{[originallemma][0][]00L8}{[1][0][]source}{}{}{}}
\makeatother
\crefname{originallemma}{original lemma}{original lemmas}
\Crefname{originallemma}{Original lemma}{Original lemmas}
\makeatletter
\newlabel{dualizing-lemma-universally-catenary}{{0A80}{source}{}{sourcepointer-18}{}}
\newlabel{dualizing-lemma-universally-catenary@cref}{{[originallemma][0][]0A80}{[1][0][]source}{}{}{}}
\makeatother
\crefname{originallemma}{original lemma}{original lemmas}
\Crefname{originallemma}{Original lemma}{Original lemmas}
\makeatletter
\newlabel{local-cohomology-lemma-catenary-S2-equidimensional}{{0FIW}{source}{}{sourcepointer-19}{}}
\newlabel{local-cohomology-lemma-catenary-S2-equidimensional@cref}{{[originallemma][0][]0FIW}{[1][0][]source}{}{}{}}
\makeatother
\title{Kollar injectivity for Picard groups of punctured spectra under a depth condition}
\author{The Stacks Project Authors}
\date{stacks-project:algebraization.tex:local-snapshot-20261011; GNU FDL 1.2 or later}
\begin{document}
\maketitle
\noindent\textit{Editorial note.} Author statement and proof are retained without a mathematical correction. An earlier assistant raised an unconfirmed concern about the Ext/image wording at source lines 8600–8659; this preparation neither adopts the proposed correction nor adjudicates correctness.\par\medskip
\noindent\textit{Editorial note.} Only the original depth/dimension hypotheses are retained. Coherent triples and their polynomial/positivity lemmas are uncounted supporting material.\par\medskip
\noindent\textit{Allowed standard prerequisites (editorial).}\begin{itemize}
\item Explicit prerequisite boundary. Local cohomology and normalized dualizing complexes, Matlis duality, depth/support inequalities, Hilbert–Samuel growth in dimension one, existence of a nonvanishing section on finitely many points of a quasi-affine scheme, and catenary S2 dimension theory. Coherent triples and their polynomial/positivity lemmas are retained local proof.
\end{itemize}\medskip
\ifdefined\hypertarget\hypertarget{sourcepointer-1}{}\fi
\noindent\textit{Source reference.} Original lemma 0BK0: focused-seeds/AGP/sources/stacks-snapshot-20261011/tags.txt, source lines 14543--14543. Stable Stacks tag for an expressly allowed foundational or notational input; the counted construction remains the two concrete core obligations. The retained local core is not replaced by this pointer. See standard\_prerequisites.\par\medskip
\ifdefined\hypertarget\hypertarget{sourcepointer-2}{}\fi
\noindent\textit{Source reference.} Original lemma 01YB: focused-seeds/AGP/sources/stacks-snapshot-20261011/tags.txt, source lines 2550--2550. Stable Stacks tag for an expressly allowed foundational or notational input; the counted construction remains the two concrete core obligations. The retained local core is not replaced by this pointer. See standard\_prerequisites.\par\medskip
\ifdefined\hypertarget\hypertarget{sourcepointer-3}{}\fi
\noindent\textit{Source reference.} Original lemma 00L5: focused-seeds/AGP/sources/stacks-snapshot-20261011/tags.txt, source lines 775--775. Stable Stacks tag for an expressly allowed foundational or notational input; the counted construction remains the two concrete core obligations. The retained local core is not replaced by this pointer. See standard\_prerequisites.\par\medskip
\ifdefined\hypertarget\hypertarget{sourcepointer-4}{}\fi
\noindent\textit{Source reference.} Original lemma 00IV: focused-seeds/AGP/sources/stacks-snapshot-20261011/tags.txt, source lines 693--693. Stable Stacks tag for an expressly allowed foundational or notational input; the counted construction remains the two concrete core obligations. The retained local core is not replaced by this pointer. See standard\_prerequisites.\par\medskip
\ifdefined\hypertarget\hypertarget{sourcepointer-5}{}\fi
\noindent\textit{Source reference.} Original lemma 00IN: focused-seeds/AGP/sources/stacks-snapshot-20261011/tags.txt, source lines 685--685. Stable Stacks tag for an expressly allowed foundational or notational input; the counted construction remains the two concrete core obligations. The retained local core is not replaced by this pointer. See standard\_prerequisites.\par\medskip
\ifdefined\hypertarget\hypertarget{sourcepointer-6}{}\fi
\noindent\textit{Source reference.} Original lemma 0F20: focused-seeds/AGP/sources/stacks-snapshot-20261011/tags.txt, source lines 18815--18815. Stable Stacks tag for an expressly allowed foundational or notational input; the counted construction remains the two concrete core obligations. The retained local core is not replaced by this pointer. See standard\_prerequisites.\par\medskip
\ifdefined\hypertarget\hypertarget{sourcepointer-7}{}\fi
\noindent\textit{Source reference.} Original lemma 00JZ: focused-seeds/AGP/sources/stacks-snapshot-20261011/tags.txt, source lines 733--733. Stable Stacks tag for an expressly allowed foundational or notational input; the counted construction remains the two concrete core obligations. The retained local core is not replaced by this pointer. See standard\_prerequisites.\par\medskip
\ifdefined\hypertarget\hypertarget{sourcepointer-8}{}\fi
\noindent\textit{Source reference.} Original lemma 090R: focused-seeds/AGP/sources/stacks-snapshot-20261011/tags.txt, source lines 11419--11419. Stable Stacks tag for an expressly allowed foundational or notational input; the counted construction remains the two concrete core obligations. The retained local core is not replaced by this pointer. See standard\_prerequisites.\par\medskip
\ifdefined\hypertarget\hypertarget{sourcepointer-9}{}\fi
\noindent\textit{Source reference.} Original lemma 0AVZ: focused-seeds/AGP/sources/stacks-snapshot-20261011/tags.txt, source lines 13702--13702. Stable Stacks tag for an expressly allowed foundational or notational input; the counted construction remains the two concrete core obligations. The retained local core is not replaced by this pointer. See standard\_prerequisites.\par\medskip
\ifdefined\hypertarget\hypertarget{sourcepointer-10}{}\fi
\noindent\textit{Source reference.} Original lemma 0BLT: focused-seeds/AGP/sources/stacks-snapshot-20261011/tags.txt, source lines 14606--14606. Stable Stacks tag for an expressly allowed foundational or notational input; the counted construction remains the two concrete core obligations. The retained local core is not replaced by this pointer. See standard\_prerequisites.\par\medskip
\ifdefined\hypertarget\hypertarget{sourcepointer-11}{}\fi
\noindent\textit{Source reference.} Original lemma 0EY0: focused-seeds/AGP/sources/stacks-snapshot-20261011/tags.txt, source lines 18675--18675. Stable Stacks tag for an expressly allowed foundational or notational input; the counted construction remains the two concrete core obligations. The retained local core is not replaced by this pointer. See standard\_prerequisites.\par\medskip
\ifdefined\hypertarget\hypertarget{sourcepointer-12}{}\fi
\noindent\textit{Source reference.} Original lemma 00LX: focused-seeds/AGP/sources/stacks-snapshot-20261011/tags.txt, source lines 803--803. Stable Stacks tag for an expressly allowed foundational or notational input; the counted construction remains the two concrete core obligations. The retained local core is not replaced by this pointer. See standard\_prerequisites.\par\medskip
\ifdefined\hypertarget\hypertarget{sourcepointer-13}{}\fi
\noindent\textit{Source reference.} Original lemma 00LK: focused-seeds/AGP/sources/stacks-snapshot-20261011/tags.txt, source lines 790--790. Stable Stacks tag for an expressly allowed foundational or notational input; the counted construction remains the two concrete core obligations. The retained local core is not replaced by this pointer. See standard\_prerequisites.\par\medskip
\ifdefined\hypertarget\hypertarget{sourcepointer-14}{}\fi
\noindent\textit{Source reference.} Original lemma 0BPY: focused-seeds/AGP/sources/stacks-snapshot-20261011/tags.txt, source lines 14716--14716. Stable Stacks tag for an expressly allowed foundational or notational input; the counted construction remains the two concrete core obligations. The retained local core is not replaced by this pointer. See standard\_prerequisites.\par\medskip
\ifdefined\hypertarget\hypertarget{sourcepointer-15}{}\fi
\noindent\textit{Source reference.} Original lemma 0AAK: focused-seeds/AGP/sources/stacks-snapshot-20261011/tags.txt, source lines 12988--12988. Stable Stacks tag for an expressly allowed foundational or notational input; the counted construction remains the two concrete core obligations. The retained local core is not replaced by this pointer. See standard\_prerequisites.\par\medskip
\ifdefined\hypertarget\hypertarget{sourcepointer-16}{}\fi
\noindent\textit{Source reference.} Original lemma 0DWZ: focused-seeds/AGP/sources/stacks-snapshot-20261011/tags.txt, source lines 17414--17414. Stable Stacks tag for an expressly allowed foundational or notational input; the counted construction remains the two concrete core obligations. The retained local core is not replaced by this pointer. See standard\_prerequisites.\par\medskip
\ifdefined\hypertarget\hypertarget{sourcepointer-17}{}\fi
\noindent\textit{Source reference.} Original lemma 00L8: focused-seeds/AGP/sources/stacks-snapshot-20261011/tags.txt, source lines 778--778. Stable Stacks tag for an expressly allowed foundational or notational input; the counted construction remains the two concrete core obligations. The retained local core is not replaced by this pointer. See standard\_prerequisites.\par\medskip
\ifdefined\hypertarget\hypertarget{sourcepointer-18}{}\fi
\noindent\textit{Source reference.} Original lemma 0A80: focused-seeds/AGP/sources/stacks-snapshot-20261011/tags.txt, source lines 12898--12898. Stable Stacks tag for an expressly allowed foundational or notational input; the counted construction remains the two concrete core obligations. The retained local core is not replaced by this pointer. See standard\_prerequisites.\par\medskip
\ifdefined\hypertarget\hypertarget{sourcepointer-19}{}\fi
\noindent\textit{Source reference.} Original lemma 0FIW: focused-seeds/AGP/sources/stacks-snapshot-20261011/tags.txt, source lines 19406--19406. Stable Stacks tag for an expressly allowed foundational or notational input; the counted construction remains the two concrete core obligations. The retained local core is not replaced by this pointer. See standard\_prerequisites.\par\medskip
\begin{proposition}[Koll\'ar]
\label{proposition-injective-pic}
\begin{reference}
\cite[Theorem 1.9]{Kollar-map-pic}
\end{reference}
Let $(A, \mathfrak m)$ be a Noetherian local ring. Let $f \in \mathfrak m$.
Assume
\begin{enumerate}
\item $A$ has a dualizing complex,
\item $f$ is a nonzerodivisor,
\item $\text{depth}(A/fA) \geq 2$, or equivalently $\text{depth}(A) \geq 3$,
\item if $f \in \mathfrak p \subset A$ is a prime ideal with
$\dim(A/\mathfrak p) = 2$, then $\text{depth}(A_\mathfrak p) \geq 2$.
\end{enumerate}
Let $U$, resp.\ $U_0$ be the punctured spectrum of $A$, resp.\ $A/fA$. The map
$$
\Pic(U) \to \Pic(U_0)
$$
is injective. Finally, if (1), (2), (3), $A$ is $(S_2)$, and
$\dim(A) \geq 4$, then (4) holds.
\end{proposition}

\section{Coherent triples}
\label{section-coherent-triples}

\noindent
Let $(A, \mathfrak m)$ be a Noetherian local ring.
Let $f \in \mathfrak m$ be a nonzerodivisor. Set
$X = \Spec(A)$, $X_0 = \Spec(A/fA)$, $U = X \setminus V(\mathfrak m)$, and
$U_0 = U \cap X_0$.
We say $(\mathcal{F}, \mathcal{F}_0, \alpha)$ is a {\it coherent triple}
if we have
\begin{enumerate}
\item $\mathcal{F}$ is a coherent $\mathcal{O}_U$-module such that
$f : \mathcal{F} \to \mathcal{F}$ is injective,
\item $\mathcal{F}_0$ is a coherent $\mathcal{O}_{X_0}$-module,
\item $\alpha : \mathcal{F}/f\mathcal{F} \to \mathcal{F}_0|_{U_0}$
is an isomorphism.
\end{enumerate}
There is an obvious notion of a {\it morphism of coherent triples}
which turns the collection of all coherent triples into a category.

\medskip\noindent
The category of coherent triples is additive but not abelian.
However, it is clear what a short exact sequence of coherent
triples is.

\medskip\noindent
Given two coherent triples $(\mathcal{F}, \mathcal{F}_0, \alpha)$
and $(\mathcal{G}, \mathcal{G}_0, \beta)$ it may not be the case that
$(\mathcal{F} \otimes_{\mathcal{O}_U} \mathcal{G},
\mathcal{F}_0  \otimes_{\mathcal{O}_{X_0}} \mathcal{G}_0,
\alpha \otimes \beta)$ is a coherent triple\footnote{Namely, it
isn't necessarily the case that $f$
is injective on $\mathcal{F} \otimes_{\mathcal{O}_U} \mathcal{G}$.}.
However, if the stalks $\mathcal{G}_x$ are
free for all $x \in U_0$, then this does hold.

\medskip\noindent
We will say the coherent triple $(\mathcal{G}, \mathcal{G}_0, \beta)$
is {\it locally free}, resp.\ {\it invertible}
if $\mathcal{G}$ and $\mathcal{G}_0$
are locally free, resp.\ invertible modules. In this case tensoring
with $(\mathcal{G}, \mathcal{G}_0, \beta)$ makes sense (see above)
and turns short exact sequences of coherent triples into short exact
sequences of coherent triples.

\begin{lemma}
\label{lemma-prepare-chi-triple}
For any coherent triple $(\mathcal{F}, \mathcal{F}_0, \alpha)$
there exists a coherent $\mathcal{O}_X$-module $\mathcal{F}'$
such that $f : \mathcal{F}' \to \mathcal{F}'$ is injective,
an isomorphism $\alpha' : \mathcal{F}'|_U \to \mathcal{F}$, and a map
$\alpha'_0 : \mathcal{F}'/f\mathcal{F}' \to \mathcal{F}_0$
such that $\alpha \circ (\alpha' \bmod f) = \alpha'_0|_{U_0}$.
\end{lemma}

\begin{proof}
Choose a finite $A$-module $M$ such that $\mathcal{F}$ is the restriction
to $U$ of the coherent $\mathcal{O}_X$-module associated to $M$, see
Local Cohomology, Lemma
\ref{local-cohomology-lemma-finiteness-pushforwards-and-H1-local}.
Since $\mathcal{F}$ is $f$-torsion free, we may replace $M$ by
its quotient by $f$-power torsion.
On the other hand, let $M_0 = \Gamma(X_0, \mathcal{F}_0)$
so that $\mathcal{F}_0$ is the coherent $\mathcal{O}_{X_0}$-module
associated to the finite $A/fA$-module $M_0$.
By Cohomology of Schemes, Lemma \ref{coherent-lemma-homs-over-open}
there exists an $n$ such that the
isomorphism $\alpha_0$ corresponds to an $A/fA$-module homomorphism
$\mathfrak m^n M/fM \to M_0$ (whose kernel and cokernel
are annihilated by a power of $\mathfrak m$, but we don't need this).
Thus if we take $M' = \mathfrak m^n M$ and we let
$\mathcal{F}'$ be the coherent $\mathcal{O}_X$-module
associated to $M'$, then the lemma is clear.
\end{proof}

\noindent
Let $(\mathcal{F}, \mathcal{F}_0, \alpha)$ be a coherent triple.
Choose $\mathcal{F}', \alpha', \alpha'_0$ as in
Lemma \ref{lemma-prepare-chi-triple}. Set
\begin{equation}
\label{equation-chi-triple}
\chi(\mathcal{F}, \mathcal{F}_0, \alpha) =
\text{length}_A(\Coker(\alpha'_0)) - 
\text{length}_A(\Ker(\alpha'_0))
\end{equation}
The expression on the right makes sense as $\alpha'_0$ is an isomorphism
over $U_0$ and hence its kernel and coherent are coherent modules supported
on $\{\mathfrak m\}$ which therefore have finite length
(Algebra, Lemma \ref{algebra-lemma-support-point}).

\begin{lemma}
\label{lemma-well-defined-chi-triple}
The quantity $\chi(\mathcal{F}, \mathcal{F}_0, \alpha)$ in
(\ref{equation-chi-triple}) does not depend on the choice of
$\mathcal{F}', \alpha', \alpha'_0$ as in Lemma \ref{lemma-prepare-chi-triple}.
\end{lemma}

\begin{proof}
Let $\mathcal{F}', \alpha', \alpha'_0$ and
$\mathcal{F}'', \alpha'', \alpha''_0$ be two such choices.
For $n > 0$ set $\mathcal{F}'_n = \mathfrak m^n \mathcal{F}'$.
By Cohomology of Schemes, Lemma \ref{coherent-lemma-homs-over-open}
for some $n$
there exists an $\mathcal{O}_X$-module map $\mathcal{F}'_n \to \mathcal{F}''$
agreeing with the identification
$\mathcal{F}''|_U = \mathcal{F}'|_U$ determined by $\alpha'$ and $\alpha''$.
Then the diagram
$$
\xymatrix{
\mathcal{F}'_n/f\mathcal{F}'_n \ar[r] \ar[d] &
\mathcal{F}'/f\mathcal{F}' \ar[d]^{\alpha_0'} \\
\mathcal{F}''/f\mathcal{F}'' \ar[r]^{\alpha_0''} &
\mathcal{F}_0
}
$$
is commutative after restricting to $U_0$. Hence by
Cohomology of Schemes, Lemma \ref{coherent-lemma-homs-over-open}
it is commutative after restricting to
$\mathfrak m^l(\mathcal{F}'_n/f\mathcal{F}'_n)$ for some $l > 0$. Since
$\mathcal{F}'_{n + l}/f\mathcal{F}'_{n + l} \to \mathcal{F}'_n/f\mathcal{F}'_n$
factors through $\mathfrak m^l(\mathcal{F}'_n/f\mathcal{F}'_n)$
we see that after replacing $n$ by $n + l$ the diagram
is commutative. In other words, we have found a third choice
$\mathcal{F}''', \alpha''', \alpha'''_0$
such that there are maps $\mathcal{F}''' \to \mathcal{F}''$
and $\mathcal{F}''' \to \mathcal{F}'$ over $X$
compatible with the maps over $U$ and $X_0$. This reduces us to
the case discussed in the next paragraph.

\medskip\noindent
Assume we have a map $\mathcal{F}'' \to \mathcal{F}'$ over $X$ compatible with
$\alpha', \alpha''$ over $U$ and with $\alpha'_0, \alpha''_0$ over $X_0$.
Observe that $\mathcal{F}'' \to \mathcal{F}'$ is injective as it is an
isomorphism over $U$ and since $f : \mathcal{F}'' \to \mathcal{F}''$
is injective. Clearly $\mathcal{F}'/\mathcal{F}''$ is supported on
$\{\mathfrak m\}$ hence has finite length. We have the maps
of coherent $\mathcal{O}_{X_0}$-modules
$$
\mathcal{F}''/f\mathcal{F}'' \to
\mathcal{F}'/f\mathcal{F}' \xrightarrow{\alpha'_0}
\mathcal{F}_0
$$
whose composition is $\alpha''_0$ and which are isomorphisms over $U_0$.
Elementary homological algebra gives a $6$-term exact sequence
$$
\begin{matrix}
0 \to
\Ker(\mathcal{F}''/f\mathcal{F}'' \to \mathcal{F}'/f\mathcal{F}') \to
\Ker(\alpha''_0) \to
\Ker(\alpha'_0) \to \\
\Coker(\mathcal{F}''/f\mathcal{F}'' \to \mathcal{F}'/f\mathcal{F}') \to
\Coker(\alpha''_0) \to
\Coker(\alpha'_0) \to 0
\end{matrix}
$$
By additivity of lengths (Algebra, Lemma \ref{algebra-lemma-length-additive})
we find that it suffices to show that
$$
\text{length}_A(
\Coker(\mathcal{F}''/f\mathcal{F}'' \to \mathcal{F}'/f\mathcal{F}')) -
\text{length}_A(
\Ker(\mathcal{F}''/f\mathcal{F}'' \to \mathcal{F}'/f\mathcal{F}')) = 0
$$
This follows from applying the snake lemma to
the diagram
$$
\xymatrix{
0 \ar[r] &
\mathcal{F}'' \ar[r]_f \ar[d] &
\mathcal{F}'' \ar[r] \ar[d] &
\mathcal{F}''/f\mathcal{F}'' \ar[r] \ar[d] &
0 \\
0 \ar[r] &
\mathcal{F}' \ar[r]^f &
\mathcal{F}' \ar[r] &
\mathcal{F}'/f\mathcal{F}' \ar[r] &
0
}
$$
and the fact that $\mathcal{F}'/\mathcal{F}''$ has finite length.
\end{proof}

\begin{lemma}
\label{lemma-ses-chi-triple}
We have
$\chi(\mathcal{G}, \mathcal{G}_0, \beta) =
\chi(\mathcal{F}, \mathcal{F}_0, \alpha) +
\chi(\mathcal{H}, \mathcal{H}_0, \gamma)$ if
$$
0 \to
(\mathcal{F}, \mathcal{F}_0, \alpha) \to
(\mathcal{G}, \mathcal{G}_0, \beta) \to
(\mathcal{H}, \mathcal{H}_0, \gamma)
\to 0
$$
is a short exact sequence of coherent triples.
\end{lemma}

\begin{proof}
Choose $\mathcal{G}', \beta', \beta'_0$ as in
Lemma \ref{lemma-prepare-chi-triple}
for the triple $(\mathcal{G}, \mathcal{G}_0, \beta)$.
Denote $j : U \to X$ the inclusion morphism.
Let $\mathcal{F}' \subset \mathcal{G}'$
be the kernel of the composition
$$
\mathcal{G}' \xrightarrow{\beta'} j_*\mathcal{G} \to j_*\mathcal{H}
$$
Observe that $\mathcal{H}' = \mathcal{G}'/\mathcal{F}'$
is a coherent subsheaf of $j_*\mathcal{H}$ and hence
$f : \mathcal{H}' \to \mathcal{H}'$ is injective.
Hence by the snake lemma we obtain a short exact sequence
$$
0 \to \mathcal{F}'/f\mathcal{F}' \to
\mathcal{G}'/f\mathcal{G}' \to
\mathcal{H}'/f\mathcal{H}' \to 0
$$
We have isomorphisms
$\alpha' : \mathcal{F}'|_U \to \mathcal{F}$,
$\beta' : \mathcal{G}'|_U \to \mathcal{G}$, and
$\gamma' : \mathcal{H}'|_U \to \mathcal{H}$ by construction.
To finish the proof we'll need to construct maps
$\alpha'_0 : \mathcal{F}'/f\mathcal{F}' \to \mathcal{F}_0$ and
$\gamma'_0 : \mathcal{H}'/f\mathcal{H}' \to \mathcal{H}_0$ as in
Lemma \ref{lemma-prepare-chi-triple} and fitting into
a commutative diagram
$$
\xymatrix{
0 \ar[r] &
\mathcal{F}'/f\mathcal{F}' \ar[r] \ar@{..>}[d]^{\alpha'_0} &
\mathcal{G}'/f\mathcal{G}' \ar[r] \ar[d]^{\beta'_0} &
\mathcal{H}'/f\mathcal{H}' \ar[r] \ar@{..>}[d]^{\gamma'_0} &
0 \\
0 \ar[r] &
\mathcal{F}_0 \ar[r] &
\mathcal{G}_0 \ar[r] &
\mathcal{H}_0 \ar[r] &
0
}
$$
However, this may not be possible with our initial choice of $\mathcal{G}'$.
From the displayed diagram we see the obstruction is
exactly the composition
$$
\delta :
\mathcal{F}'/f\mathcal{F}' \to
\mathcal{G}'/f\mathcal{G}' \xrightarrow{\beta'_0}
\mathcal{G}_0 \to
\mathcal{H}_0
$$
Note that the restriction of $\delta$ to $U_0$ is zero by our choice of
$\mathcal{F}'$ and $\mathcal{H}'$. Hence by
Cohomology of Schemes, Lemma \ref{coherent-lemma-homs-over-open}
there exists an $k > 0$ such that
$\delta$ vanishes on $\mathfrak m^k \cdot (\mathcal{F}'/f\mathcal{F}')$.
For $n > k$ set $\mathcal{G}'_n = \mathfrak m^n \mathcal{G}'$,
$\mathcal{F}'_n = \mathcal{G}'_n \cap \mathcal{F}'$, and
$\mathcal{H}'_n = \mathcal{G}'_n/\mathcal{F}'_n$.
Observe that $\beta'_0$ can be composed with
$\mathcal{G}'_n/f\mathcal{G}'_n \to \mathcal{G}'/f\mathcal{G}'$
to give a map
$\beta'_{n, 0} : \mathcal{G}'_n/f\mathcal{G}'_n \to \mathcal{G}_0$
as in Lemma \ref{lemma-prepare-chi-triple}.
By Artin-Rees (Algebra, Lemma \ref{algebra-lemma-Artin-Rees})
we may choose $n$ such that
$\mathcal{F}'_n \subset \mathfrak m^k \mathcal{F}'$.
As above the maps
$f : \mathcal{F}'_n \to \mathcal{F}'_n$,
$f : \mathcal{G}'_n \to \mathcal{G}'_n$, and
$f : \mathcal{H}'_n \to \mathcal{H}'_n$ are injective
and as above using the snake lemma we obtain a short exact
sequence
$$
0 \to \mathcal{F}'_n/f\mathcal{F}'_n \to
\mathcal{G}'_n/f\mathcal{G}'_n \to
\mathcal{H}'_n/f\mathcal{H}'_n \to 0
$$
As above we have isomorphisms
$\alpha'_n : \mathcal{F}'_n|_U \to \mathcal{F}$,
$\beta'_n : \mathcal{G}'_n|_U \to \mathcal{G}$, and
$\gamma'_n : \mathcal{H}'_n|_U \to \mathcal{H}$.
We consider the obstruction
$$
\delta_n :
\mathcal{F}'_n/f\mathcal{F}'_n \to
\mathcal{G}'_n/f\mathcal{G}'_n
\xrightarrow{\beta'_{n, 0}}
\mathcal{G}_0 \to
\mathcal{H}_0
$$
as before. However, the commutative diagram
$$
\xymatrix{
\mathcal{F}'_n/f\mathcal{F}'_n \ar[r] \ar[d] &
\mathcal{G}'_n/f\mathcal{G}'_n \ar[r]_{\beta'_{n, 0}} \ar[d] &
\mathcal{G}_0 \ar[r] \ar[d] &
\mathcal{H}_0 \ar[d] \\
\mathcal{F}'/f\mathcal{F}' \ar[r] &
\mathcal{G}'/f\mathcal{G}' \ar[r]^{\beta'_0} &
\mathcal{G}_0 \ar[r] &
\mathcal{H}_0
}
$$
our choice of $n$ and our observation about $\delta$
show that $\delta_n = 0$.
This produces the desired maps
$\alpha'_{n, 0} : \mathcal{F}'_n/f\mathcal{F}'_n \to \mathcal{F}_0$, and
$\gamma'_{n, 0} : \mathcal{H}'_n/f\mathcal{H}'_n \to \mathcal{H}_0$.
OK, so we may use
$\mathcal{F}'_n, \alpha'_n, \alpha'_{n, 0}$,
$\mathcal{G}'_n, \beta'_n, \beta'_{n, 0}$, and
$\mathcal{H}'_n, \gamma'_n, \gamma'_{n, 0}$
to compute
$\chi(\mathcal{F}, \mathcal{F}_0, \alpha)$,
$\chi(\mathcal{G}, \mathcal{G}_0, \beta)$, and
$\chi(\mathcal{H}, \mathcal{H}_0, \gamma)$.
Now finally the lemma follows from
an application of the snake lemma to
$$
\xymatrix{
0 \ar[r] &
\mathcal{F}'_n/f\mathcal{F}'_n \ar[r] \ar[d] &
\mathcal{G}'_n/f\mathcal{G}'_n \ar[r] \ar[d] &
\mathcal{H}'_n/f\mathcal{H}'_n \ar[r] \ar[d] &
0 \\
0 \ar[r] &
\mathcal{F}_0 \ar[r] &
\mathcal{G}_0 \ar[r] &
\mathcal{H}_0 \ar[r] &
0
}
$$
and additivity of lengths (Algebra, Lemma \ref{algebra-lemma-length-additive}).
\end{proof}

\begin{proposition}
\label{proposition-hilbert-triple}
Let $(\mathcal{F}, \mathcal{F}_0, \alpha)$ be a coherent triple.
Let $(\mathcal{L}, \mathcal{L}_0, \lambda)$ be an invertible coherent
triple. Then the function
$$
\mathbf{Z} \longrightarrow \mathbf{Z},\quad
n \longmapsto
\chi((\mathcal{F}, \mathcal{F}_0, \alpha) \otimes
(\mathcal{L}, \mathcal{L}_0, \lambda)^{\otimes n})
$$
is a polynomial of degree $\leq \dim(\text{Supp}(\mathcal{F}))$.
\end{proposition}

\noindent
More precisely, if $\mathcal{F} = 0$, then the function is constant.
If $\mathcal{F}$ has finite support in $U$, then the function is constant.
If the support of $\mathcal{F}$ in $U$ has dimension $1$, i.e., the
closure of the support of $\mathcal{F}$ in $X$ has dimension $2$, then
the function is linear, etc.

\begin{proof}
We will prove this by induction on the dimension of the support of
$\mathcal{F}$.

\medskip\noindent
The base case is when $\mathcal{F} = 0$. Then either
$\mathcal{F}_0$ is zero or its support is $\{\mathfrak m\}$.
In this case we have
$$
(\mathcal{F}, \mathcal{F}_0, \alpha) \otimes
(\mathcal{L}, \mathcal{L}_0, \lambda)^{\otimes n} =
(0, \mathcal{F}_0 \otimes \mathcal{L}_0^{\otimes n}, 0) \cong
(0, \mathcal{F}_0, 0)
$$
Thus the function of the lemma is constant with value equal
to the length of $\mathcal{F}_0$.

\medskip\noindent
Induction step. Assume the support of $\mathcal{F}$ is nonempty.
Let $\mathcal{G}_0 \subset \mathcal{F}_0$ denote the submodule
of sections supported on $\{\mathfrak m\}$. Then we get a short
exact sequence
$$
0 \to (0, \mathcal{G}_0, 0) \to
(\mathcal{F}, \mathcal{F}_0, \alpha) \to
(\mathcal{F}, \mathcal{F}_0/\mathcal{G}_0, \alpha) \to 0
$$
This sequence remains exact if we tensor by the invertible
coherent triple $(\mathcal{L}, \mathcal{L}_0, \lambda)$, see
discussion above. Thus by additivity of $\chi$
(Lemma \ref{lemma-ses-chi-triple})
and the base case explained above, it suffices to prove
the induction step for
$(\mathcal{F}, \mathcal{F}_0/\mathcal{G}_0, \alpha)$.
In this way we see that we may assume $\mathfrak m$ is not
an associated point of $\mathcal{F}_0$.

\medskip\noindent
Let $T = \text{Ass}(\mathcal{F}) \cup \text{Ass}(\mathcal{F}/f\mathcal{F})$.
Since $U$ is quasi-affine, we can find $s \in \Gamma(U, \mathcal{L})$
which does not vanish at any $u \in T$, see
Properties, Lemma
\ref{properties-lemma-quasi-affine-invertible-nonvanishing-section}.
After multiplying $s$ by a suitable element of $\mathfrak m$
we may assume $\lambda(s \bmod f) = s_0|_{U_0}$ for some
$s_0 \in \Gamma(X_0, \mathcal{L}_0)$; details omitted.
We obtain a morphism
$$
(s, s_0) :
(\mathcal{O}_U, \mathcal{O}_{X_0}, 1)
\longrightarrow
(\mathcal{L}, \mathcal{L}_0, \lambda)
$$
in the category of coherent triples. Let
$\mathcal{G} = \Coker(s : \mathcal{F} \to \mathcal{F} \otimes \mathcal{L})$
and
$\mathcal{G}_0 = \Coker(s_0 : \mathcal{F}_0 \to
\mathcal{F}_0 \otimes \mathcal{L}_0)$. Observe that $s_0 : \mathcal{F}_0 \to
\mathcal{F}_0 \otimes \mathcal{L}_0$ is injective as it is injective
on $U_0$ by our choice of $s$ and as $\mathfrak m$ isn't an
associated point of $\mathcal{F}_0$. It follows that
there exists an
isomorphism $\beta : \mathcal{G}/f\mathcal{G} \to \mathcal{G}_0|_{U_0}$
such that we obtain a short exact sequence
$$
0 \to
(\mathcal{F}, \mathcal{F}_0, \alpha) \to
(\mathcal{F}, \mathcal{F}_0, \alpha) \otimes
(\mathcal{L}, \mathcal{L}_0, \lambda) \to
(\mathcal{G}, \mathcal{G}_0, \beta) \to 0
$$
By induction on the dimension of the support we know the proposition
holds for the coherent triple $(\mathcal{G}, \mathcal{G}_0, \beta)$.
Using the additivity of Lemma \ref{lemma-ses-chi-triple}
we see that
$$
n \longmapsto 
\chi((\mathcal{F}, \mathcal{F}_0, \alpha) \otimes
(\mathcal{L}, \mathcal{L}_0, \lambda)^{\otimes n + 1})
-
\chi((\mathcal{F}, \mathcal{F}_0, \alpha) \otimes
(\mathcal{L}, \mathcal{L}_0, \lambda)^{\otimes n})
$$
is a polynomial. We conclude by a variant of
Algebra, Lemma \ref{algebra-lemma-numerical-polynomial}
for functions defined for all integers (details omitted).
\end{proof}

\begin{lemma}
\label{lemma-nonnegative-chi-triple}
Assume $\text{depth}(A) \geq 3$ or equivalently
$\text{depth}(A/fA) \geq 2$. Let $(\mathcal{L}, \mathcal{L}_0, \lambda)$
be an invertible coherent triple. Then
$$
\chi(\mathcal{L}, \mathcal{L}_0, \lambda) =
\text{length}_A \Coker(\Gamma(U, \mathcal{L}) \to \Gamma(U_0, \mathcal{L}_0))
$$
and in particular this is $\geq 0$. Moreover, 
$\chi(\mathcal{L}, \mathcal{L}_0, \lambda) = 0$ if and only if
$\mathcal{L} \cong \mathcal{O}_U$.
\end{lemma}

\begin{proof}
The equivalence of the depth conditions follows from
Algebra, Lemma \ref{algebra-lemma-depth-drops-by-one}.
By the depth condition we see that
$\Gamma(U, \mathcal{O}_U) = A$ and
$\Gamma(U_0, \mathcal{O}_{U_0}) = A/fA$, see
Dualizing Complexes, Lemma \ref{dualizing-lemma-depth} and
Local Cohomology, Lemma
\ref{local-cohomology-lemma-finiteness-pushforwards-and-H1-local}.
Using Local Cohomology, Lemma
\ref{local-cohomology-lemma-finiteness-for-finite-locally-free}
we find that $M = \Gamma(U, \mathcal{L})$ is a finite $A$-module.
This in turn implies $\text{depth}(M) \geq 2$ for example by
part (4) of Local Cohomology, Lemma
\ref{local-cohomology-lemma-finiteness-pushforwards-and-H1-local}
or by Divisors, Lemma \ref{divisors-lemma-depth-pushforward}.
Also, we have $\mathcal{L}_0 \cong \mathcal{O}_{X_0}$
as $X_0$ is a local scheme. Hence we also see that
$M_0 = \Gamma(X_0, \mathcal{L}_0) = \Gamma(U_0, \mathcal{L}_0|_{U_0})$
and that this module is isomorphic to $A/fA$.

\medskip\noindent
By the above $\mathcal{F}' = \widetilde{M}$ is a coherent
$\mathcal{O}_X$-module whose restriction to $U$ is isomorphic to $\mathcal{L}$.
The isomorphism $\lambda : \mathcal{L}/f\mathcal{L} \to \mathcal{L}_0|_{U_0}$
determines a map $M/fM \to M_0$ on global sections
which is an isomorphism over $U_0$.
Since $\text{depth}(M) \geq 2$ we see
that $H^0_\mathfrak m(M/fM) = 0$ and it follows that
$M/fM \to M_0$ is injective. Thus by definition
$$
\chi(\mathcal{L}, \mathcal{L}_0, \lambda) =
\text{length}_A \Coker(M/fM \to M_0)
$$
which gives the first statement of the lemma.

\medskip\noindent
Finally, if this length is $0$, then $M \to M_0$ is surjective.
Hence we can find $s \in M = \Gamma(U, \mathcal{L})$
mapping to a trivializing section of $\mathcal{L}_0$.
Consider the finite $A$-modules $K$, $Q$ defined by the exact
sequence
$$
0 \to K \to A \xrightarrow{s} M \to Q \to 0
$$
The supports of $K$ and $Q$ do not meet $U_0$ because $s$
is nonzero at points of $U_0$. Using
Algebra, Lemma \ref{algebra-lemma-depth-in-ses}
we see that $\text{depth}(K) \geq 2$ (observe that
$As \subset M$ has $\text{depth} \geq 1$ as a submodule of $M$).
Thus the support of $K$ if nonempty has dimension $\geq 2$ by
Algebra, Lemma \ref{algebra-lemma-bound-depth}.
This contradicts $\text{Supp}(M) \cap V(f) \subset \{\mathfrak m\}$
unless $K = 0$. When $K = 0$ we find that
$\text{depth}(Q) \geq 2$ and we conclude
$Q = 0$ as before. Hence $A \cong M$ and
$\mathcal{L}$ is trivial.
\end{proof}







\section{Invertible modules on punctured spectra, I}
\label{section-local-lefschetz-for-pic}

\noindent
In this section we prove some local Lefschetz theorems for the Picard group.
Some of the ideas are taken from
\cite{Kollar-pic}, \cite{Bhatt-local}, and \cite{Kollar-map-pic}.

\begin{lemma}
\label{lemma-injective-torsion-in-pic}
Let $(A, \mathfrak m)$ be a Noetherian local ring. Let $f \in \mathfrak m$
be a nonzerodivisor and assume that $\text{depth}(A/fA) \geq 2$, or equivalently
$\text{depth}(A) \geq 3$. Let $U$, resp.\ $U_0$ be the punctured
spectrum of $A$, resp.\ $A/fA$. The map
$$
\Pic(U) \to \Pic(U_0)
$$
is injective on torsion.
\end{lemma}

\begin{proof}
Let $\mathcal{L}$ be an invertible $\mathcal{O}_U$-module.
Observe that $\mathcal{L}$ maps to $0$ in $\Pic(U_0)$
if and only if we can extend $\mathcal{L}$ to an invertible
coherent triple $(\mathcal{L}, \mathcal{L}_0, \lambda)$
as in Section \ref{section-coherent-triples}.
By Proposition \ref{proposition-hilbert-triple}
the function
$$
n \longmapsto \chi((\mathcal{L}, \mathcal{L}_0, \lambda)^{\otimes n})
$$
is a polynomial. By Lemma \ref{lemma-nonnegative-chi-triple}
the value of this polynomial is zero if and only if
$\mathcal{L}^{\otimes n}$ is trivial.
Thus if $\mathcal{L}$ is torsion, then this
polynomial has infinitely many zeros, hence is
identically zero, hence $\mathcal{L}$ is trivial.
\end{proof}

\begin{proof}
Let $\mathcal{L}$ be an invertible $\mathcal{O}_U$-module.
Observe that $\mathcal{L}$ maps to $0$ in $\Pic(U_0)$
if and only if we can extend $\mathcal{L}$ to an invertible
coherent triple $(\mathcal{L}, \mathcal{L}_0, \lambda)$
as in Section \ref{section-coherent-triples}.
By Proposition \ref{proposition-hilbert-triple}
the function
$$
n \longmapsto \chi((\mathcal{L}, \mathcal{L}_0, \lambda)^{\otimes n})
$$
is a polynomial $P$. By Lemma \ref{lemma-nonnegative-chi-triple}
we have
$P(n) \geq 0$ for all $n \in \mathbf{Z}$ with equality if and only if
$\mathcal{L}^{\otimes n}$ is trivial. In particular $P(0) = 0$
and $P$ is either identically zero and we win or $P$ has even degree $\geq 2$.

\medskip\noindent
Set $M = \Gamma(U, \mathcal{L})$ and
$M_0 = \Gamma(X_0, \mathcal{L}_0) = \Gamma(U_0, \mathcal{L}_0)$.
Then $M$ is a finite $A$-module of depth $\geq 2$
and $M_0 \cong A/fA$, see proof of Lemma \ref{lemma-nonnegative-chi-triple}.
Note that $H^2_\mathfrak m(M)$ is finite $A$-module by
Local Cohomology, Lemma
\ref{local-cohomology-lemma-local-finiteness-for-finite-locally-free}
and the fact that $H^i_\mathfrak m(A) = 0$ for $i = 0, 1, 2$
since $\text{depth}(A) \geq 3$.
Consider the short exact sequence
$$
0 \to M/fM \to M_0 \to Q \to 0
$$
Lemma \ref{lemma-nonnegative-chi-triple} tells us $Q$ has finite length
equal to $\chi(\mathcal{L}, \mathcal{L}_0, \lambda)$.
We obtain $Q = H^1_\mathfrak m(M/fM)$ and
$H^i_\mathfrak m(M/fM) = H^i_\mathfrak m(M_0) \cong H^i_\mathfrak m(A/fA)$
for $i > 1$ from the long exact sequence of local cohomology
associated to the displayed short exact sequence. Consider the long
exact sequence of local cohomology associated to the sequence
$0 \to M \to M \to M/fM \to 0$. It starts with
$$
0 \to Q \to H^2_\mathfrak m(M) \to H^2_\mathfrak m(M) \to
H^2_\mathfrak m(A/fA)
$$
Using additivity of lengths we see that
$\chi(\mathcal{L}, \mathcal{L}_0, \lambda)$
is equal to the length of the image of
$H^2_\mathfrak m(M) \to H^2_\mathfrak m(A/fA)$.

\medskip\noindent
Let prove the lemma in a special case to elucidate the rest of the proof.
Namely, assume for a moment that $H^2_\mathfrak m(A/fA)$ is
a finite length module. Then
we would have $P(1) \leq \text{length}_A H^2_\mathfrak m(A/fA)$.
The exact same argument applied to $\mathcal{L}^{\otimes n}$ shows that
$P(n) \leq \text{length}_A H^2_\mathfrak m(A/fA)$ for all $n$.
Thus $P$ cannot have positive degree and we win.
In the rest of the proof we will modify this argument to give
a linear upper bound for $P(n)$ which suffices.

\medskip\noindent
Let us study the map
$H^2_\mathfrak m(M) \to H^2_\mathfrak m(M_0) \cong H^2_\mathfrak m(A/fA)$.
Choose a normalized dualizing complex $\omega_A^\bullet$ for $A$.
By local duality
(Dualizing Complexes, Lemma \ref{dualizing-lemma-special-case-local-duality})
this map is Matlis dual to the map
$$
\text{Ext}^{-2}_A(M, \omega_A^\bullet)
\longleftarrow
\text{Ext}^{-2}_A(M_0, \omega_A^\bullet)
$$
whose image therefore has the same (finite) length.
The support (if nonempty) of the finite $A$-module
$\text{Ext}^{-2}_A(M_0, \omega_A^\bullet)$ consists of
$\mathfrak m$ and a finite number of primes
$\mathfrak p_1, \ldots, \mathfrak p_r$ containing $f$ with
$\dim(A/\mathfrak p_i) = 1$. Namely, by
Local Cohomology, Lemma \ref{local-cohomology-lemma-sitting-in-degrees}
the support is contained in the set of primes $\mathfrak p \subset A$ with
$\text{depth}_{A_\mathfrak p}(M_{0, \mathfrak p}) + \dim(A/\mathfrak p) \leq 2$.
Thus it suffices to show there is no prime $\mathfrak p$ containing $f$ with
$\dim(A/\mathfrak p) = 2$ and
$\text{depth}_{A_\mathfrak p}(M_{0, \mathfrak p}) = 0$.
However, because $M_{0, \mathfrak p} \cong (A/fA)_\mathfrak p$
this would give $\text{depth}(A_\mathfrak p) = 1$ which contradicts
assumption (4).
Choose a section $t \in \Gamma(U, \mathcal{L}^{\otimes -1})$
which does not vanish in the points $\mathfrak p_1, \ldots, \mathfrak p_r$, see
Properties, Lemma
\ref{properties-lemma-quasi-affine-invertible-nonvanishing-section}.
Multiplication by $t$ on global sections determines a map $t : M \to A$
which defines an isomorphism $M_{\mathfrak p_i} \to A_{\mathfrak p_i}$ for
$i = 1, \ldots, r$. Denote $t_0 = t|_{U_0}$ the corresponding section
of $\Gamma(U_0, \mathcal{L}_0^{\otimes -1})$ which similarly determines
a map $t_0 : M_0 \to A/fA$ compatible with $t$.
We conclude that there is a commutative diagram
$$
\xymatrix{
\text{Ext}^{-2}_A(M, \omega_A^\bullet) &
\text{Ext}^{-2}_A(M_0, \omega_A^\bullet) \ar[l] \\
\text{Ext}^{-2}_A(A, \omega_A^\bullet) \ar[u]^t &
\text{Ext}^{-2}_A(A/fA, \omega_A^\bullet) \ar[l] \ar[u]_{t_0}
}
$$
It follows that the length of the image of the top horizontal
map is at most the length of $\text{Ext}^{-2}_A(A/fA, \omega_A^\bullet)$
plus the length of the cokernel of $t_0$.

\medskip\noindent
However, if we replace $\mathcal{L}$ by $\mathcal{L}^n$ for $n > 1$,
then we can use
$$
t^n :
M_n = \Gamma(U, \mathcal{L}^{\otimes n})
\longrightarrow
\Gamma(U, \mathcal{O}_U) = A
$$
instead of $t$. This replaces $t_0$ by its $n$th power.
Thus the length of the image of the map
$\text{Ext}^{-2}_A(M_n, \omega_A^\bullet) \leftarrow
\text{Ext}^{-2}_A(M_{n, 0}, \omega_A^\bullet)$
is at most the length of $\text{Ext}^{-2}_A(A/fA, \omega_A^\bullet)$
plus the length of the cokernel of
$$
t_0^n : 
\text{Ext}^{-2}_A(A/fA, \omega_A^\bullet)
\longrightarrow
\text{Ext}^{-2}_A(M_{n, 0}, \omega_A^\bullet)
$$
Via the isomorphism $M_0 \cong A/fA$ the map $t_0$ becomes
$g : A/fA \to A/fA$ for some $g \in A/fA$ and via the corresponding
isomorphisms $M_{n, 0} \cong A/fA$ the map $t_0^n$ becomes
$g^n : A/fA \to A/fA$. Thus the length of the cokernel above
is the length of the quotient of
$\text{Ext}^{-2}_A(A/fA, \omega_A^\bullet)$ by $g^n$.
Since $\text{Ext}^{-2}_A(A/fA, \omega_A^\bullet)$ is a finite $A$-module
with support $T$ of dimension $1$ and since $V(g) \cap T$
consists of the closed point by our choice of $t$
this length grows linearly in $n$ by
Algebra, Lemma \ref{algebra-lemma-support-dimension-d}.

\medskip\noindent
To finish the proof we prove the final assertion. Assume
$f \in \mathfrak m \subset A$ satisfies
(1), (2), (3), $A$ is $(S_2)$, and $\dim(A) \geq 4$.
Condition (1) implies $A$ is catenary, see
Dualizing Complexes, Lemma \ref{dualizing-lemma-universally-catenary}.
Then $\Spec(A)$ is equidimensional by Local Cohomology, Lemma
\ref{local-cohomology-lemma-catenary-S2-equidimensional}.
Thus $\dim(A_\mathfrak p) + \dim(A/\mathfrak p) \geq 4$
for every prime $\mathfrak p$ of $A$. Then
$\text{depth}(A_\mathfrak p) \geq \min(2, \dim(A_\mathfrak p))
\geq \min(2, 4 - \dim(A/\mathfrak p))$ and hence (4) holds.
\end{proof}

\bibliographystyle{plain}
\bibliography{my}
\end{document}
